\section{مقدمة}

تُعد النمذجة الرياضية للذراع الروبوتية من أهم المراحل الأساسية في تصميم وتحليل أنظمة الروبوتات، حيث تتيح فهماً دقيقاً لكيفية حركة الذراع واستجابتها للأوامر المختلفة. يعتمد الأداء الفعّال لأي نظام روبوتي على القدرة على توصيف حركته بشكل رياضي وربط متغيرات المفاصل بموقع واتجاه أداة النهاية (End-Effector).

في هذا الفصل، يتم تطوير نموذج شامل للذراع الروبوتية المستخدمة في هذا المشروع، وذلك من خلال دراسة الحركيات (Kinematics) والديناميك (Dynamics).

تشمل الحركيات جزئين رئيسيين:
\begin{itemize}
    \item \textbf{الحركيات الأمامية (\LR{Forward Kinematics}):} والتي تهدف إلى تحديد موضع واتجاه أداة النهاية اعتماداً على زوايا المفاصل.
    
    \item \textbf{الحركيات العكسية (\LR{Inverse Kinematics}):} والتي تهدف إلى حساب زوايا المفاصل المطلوبة للوصول إلى موضع واتجاه معين لأداة النهاية.
\end{itemize}

أما الديناميك، فيتعلق بدراسة القوى والعزوم المؤثرة على الذراع الروبوتية اللازمة لإحداث الحركة المطلوبة، ويتم ذلك باستخدام نماذج فيزيائية مثل طريقة لاغرانج أو نيوتن-أويلر.

تتكون الذراع الروبوتية في هذا المشروع من سلسلة من المفاصل الدورانية التي يتم التحكم بها بواسطة محركات سيرفو، مما يشكل نظاماً تسلسلياً (\LR{Serial Manipulator}). ولتوصيف هذا النظام بشكل منهجي، سيتم استخدام طريقة دينافيت-هارتنبرغ (Denavit-Hartenberg) لوصف العلاقات الهندسية بين أجزاء الذراع.

يمثل هذا الفصل الأساس النظري الذي يُبنى عليه تصميم نظام التحكم وتنفيذ الذراع الروبوتية المعتمدة على الأوامر الصوتية.

